\documentclass[12pt,a4paper]{article}

% ── Packages ──────────────────────────────────────────────────────────────────
\usepackage[utf8]{inputenc}
\usepackage[T1]{fontenc}
\usepackage[english]{babel}
\usepackage{geometry}
\usepackage{graphicx}
\usepackage{booktabs}
\usepackage{array}
\usepackage{longtable}
\usepackage{hyperref}
\usepackage{xcolor}
\usepackage{listings}
\usepackage{enumitem}
\usepackage{titlesec}
\usepackage{fancyhdr}
\usepackage{caption}
\usepackage{subcaption}
\usepackage{float}
\usepackage{parskip}
\usepackage{tcolorbox}
\tcbuselibrary{skins,breakable}

% ── Page geometry ─────────────────────────────────────────────────────────────
\geometry{top=2.5cm, bottom=2.5cm, left=2.5cm, right=2.5cm, headheight=14pt}

% ── Colors ────────────────────────────────────────────────────────────────────
\definecolor{darkblue}{RGB}{0,51,102}
\definecolor{codegreen}{rgb}{0.0,0.45,0.0}
\definecolor{codegray}{rgb}{0.5,0.5,0.5}
\definecolor{lightgray}{RGB}{248,248,248}
\definecolor{alertred}{RGB}{180,30,30}
\definecolor{infobox}{RGB}{230,240,255}
\definecolor{successgreen}{RGB}{20,120,20}

% ── tcolorbox styles ──────────────────────────────────────────────────────────
\newtcolorbox{infoblock}[1]{
  colback=infobox, colframe=darkblue, fonttitle=\bfseries,
  title=#1, breakable, left=6pt, right=6pt, top=4pt, bottom=4pt
}
\newtcolorbox{alertblock}[1]{
  colback=red!8!white, colframe=alertred, fonttitle=\bfseries\color{alertred},
  title=#1, breakable, left=6pt, right=6pt, top=4pt, bottom=4pt
}
\newtcolorbox{resultblock}[1]{
  colback=green!6!white, colframe=successgreen, fonttitle=\bfseries\color{successgreen},
  title=#1, breakable, left=6pt, right=6pt, top=4pt, bottom=4pt
}

% ── Section formatting ────────────────────────────────────────────────────────
\titleformat{\section}{\color{darkblue}\large\bfseries}{\thesection.}{0.5em}{}[\titlerule]
\titleformat{\subsection}{\color{darkblue}\normalsize\bfseries}{\thesubsection}{0.5em}{}
\titleformat{\subsubsection}{\normalsize\bfseries\itshape}{}{0em}{}

% ── Header / Footer ───────────────────────────────────────────────────────────
\pagestyle{fancy}
\fancyhf{}
\fancyhead[L]{\small\textcolor{darkblue}{IoT Networking Stack Simulation -- Track B}}
\fancyhead[R]{\small\textcolor{darkblue}{Esma S\"umer}}
\fancyfoot[C]{\small\thepage}
\renewcommand{\headrulewidth}{0.4pt}

% ── Listings ──────────────────────────────────────────────────────────────────
\lstdefinelanguage{Lua}{
  keywords={function,end,local,if,then,else,elseif,while,do,for,in,return,true,false,nil,not,and,or},
  sensitive=true,
  comment=[l]{--},
  morecomment=[s]{--[[}{]]},
  morestring=[b]",
  morestring=[b]',
}
\lstset{
  language=Lua,
  backgroundcolor=\color{lightgray},
  basicstyle=\ttfamily\footnotesize,
  keywordstyle=\color{darkblue}\bfseries,
  commentstyle=\color{codegreen}\itshape,
  stringstyle=\color{red!70!black},
  breaklines=true,
  frame=single,
  rulecolor=\color{gray!40},
  captionpos=b,
  numbers=left,
  numberstyle=\tiny\color{codegray},
  numbersep=8pt,
  showstringspaces=false,
  tabsize=2
}

% ── Hyperref ──────────────────────────────────────────────────────────────────
\hypersetup{
  colorlinks=true, linkcolor=darkblue, urlcolor=darkblue, citecolor=darkblue,
  pdftitle={IoT Networking Stack Simulation Report},
  pdfauthor={Esma Sumer}
}

% ══════════════════════════════════════════════════════════════════════════════
\begin{document}

% ── Title page ────────────────────────────────────────────────────────────────
\begin{titlepage}
  \centering
  \vspace*{1.2cm}
  {\Huge\bfseries\color{darkblue} IoT Networking Stack\\[0.4em] Simulation Report}
  \vspace{0.5cm}
  \rule{0.75\linewidth}{1.2pt}
  \vspace{0.4cm}

  {\large\itshape Track B: Implementing Book Examples in Netualizer}

  \vspace{2cm}

  {\Large\textbf{Esma S\"umer}}\\[0.4em]
  {\normalsize Computer Engineering\\Istanbul K\"ult\"ur University}

  \vspace{1.2cm}
  {\normalsize
    \textbf{Course:} COM0453 -- Internet of Things\\[0.2em]
    \textbf{Instructor:} Mehmet F.\ Yuce\\[0.2em]
    \textbf{Date:} June 2026
  }

  \vfill
  \rule{0.75\linewidth}{0.5pt}\\[0.3em]
  {\small\textcolor{gray}{Netualizer Simulation Framework~$\cdot$~Spring 2026}}
\end{titlepage}

\tableofcontents
\newpage

% ══════════════════════════════════════════════════════════════════════════════
\begin{abstract}
\noindent
This report documents the design, implementation, and analysis of seven IoT
networking protocol scenarios executed within the \textbf{Netualizer}
simulation environment as part of the COM0453 Internet of Things course
(Track~B). Five foundational protocols --- TCP, MQTT, CoAP, UDP, and mDNS ---
are implemented, each documented with its theoretical background, Netualizer
architecture diagram, Lua source code walkthrough, console output, and
engineering analysis. Two additional scenarios cover data logging with
high-resolution timestamps and deliberate multi-layer failure injection.
A comparative protocol matrix and a security threat overview conclude the
report. All experiments are fully reproducible within Netualizer without
external hardware.
\end{abstract}

\newpage

% ══════════════════════════════════════════════════════════════════════════════
\section{Introduction}
\label{sec:intro}

The explosive growth of IoT demands networking architectures fundamentally
different from conventional environments. IoT nodes --- characterised as
\emph{Constrained Nodes} in RFC~7228~\cite{RFC7228} --- operate under severe
memory, processing, and energy budgets while communicating over low-power,
lossy links. Selecting the right protocol for each layer is therefore a
critical engineering decision.

\textbf{Netualizer} provides an event-driven simulation platform that models
the complete IoT stack from the physical layer (PHY) through the application
layer. Virtual devices are scripted in Lua; the network stack is declared as
a layer chain (e.g.\ \texttt{phy -> ethernet -> ip -> tcp -> raw}). This
allows engineers to evaluate protocol behaviour, measure performance, and
inject faults without deploying physical hardware.

\subsection*{Project Scope and Objectives}

\begin{infoblock}{Track B Requirements}
Each implemented example must include: setup, run procedure, expected output,
and short explanation. Examples already implemented in class must not be
reused. Each group member implements at least 5 examples.
\end{infoblock}

This project implements and analyses the following scenarios:

\begin{enumerate}
  \item \textbf{TCP Client} --- DNS resolution + connection-oriented HTTP~GET
  \item \textbf{MQTT Sensor} --- broker-based publish/subscribe telemetry loop
  \item \textbf{CoAP Client} --- constrained REST over UDP
  \item \textbf{UDP Sender} --- raw connectionless datagram dispatch
  \item \textbf{mDNS Responder} --- zero-configuration local service discovery
  \item \textbf{Data Logging} --- timestamp-based event audit trail
  \item \textbf{Failure Injection} --- multi-layer fault simulation and analysis
\end{enumerate}

% ══════════════════════════════════════════════════════════════════════════════
\section{Simulation Environment}
\label{sec:env}

\begin{table}[H]
\centering
\caption{Simulation environment summary}
\label{tab:env}
\begin{tabular}{ll}
\toprule
\textbf{Parameter} & \textbf{Value} \\
\midrule
Simulation tool     & Netualizer (l7tr.com) \\
Scripting language  & Lua 5.x \\
Host OS             & Windows 11, 64-bit \\
Virtual subnet      & 192.168.21.0/24 \\
Default gateway     & 192.168.21.1 \\
Physical hardware   & None (fully simulated) \\
Protocol coverage   & TCP, MQTT, CoAP, UDP, mDNS \\
\bottomrule
\end{tabular}
\end{table}

Each Netualizer project consists of:
\begin{itemize}[nosep]
  \item A \textbf{config file} (.cfg) defining the virtual network stack layers
  \item A \textbf{Lua script} implementing device logic and automation
  \item The \textbf{Agents panel} showing connected virtual devices
  \item The \textbf{Events panel} logging all network state transitions with timestamps
  \item The \textbf{Output panel} displaying script console output
\end{itemize}

% ══════════════════════════════════════════════════════════════════════════════
\section{Simulation Case Studies}
\label{sec:cases}

% ── 3.1 TCP ──────────────────────────────────────────────────────────────────
\subsection{Example 1: TCP Client (Connection-Oriented HTTP GET)}
\label{sec:tcp}

\subsubsection*{Theoretical Background}

The \textbf{Transmission Control Protocol (TCP)} operates at OSI Layer~4 and
provides reliable, ordered, error-checked byte-stream delivery. Before any
data is exchanged, TCP mandates a \emph{three-way handshake}:

\begin{center}
  \texttt{SYN} $\longrightarrow$ \texttt{SYN-ACK} $\longleftarrow$ \texttt{ACK}
\end{center}

Although TCP's minimum 20-byte header and connection setup overhead are costly
for constrained devices, TCP is essential for:
\begin{itemize}[nosep]
  \item HTTP/HTTPS REST APIs
  \item Firmware-over-the-air (FOTA) updates
  \item Large payload transfers requiring guaranteed delivery
\end{itemize}

DNS resolution precedes the TCP handshake: the client queries a resolver to
map a human-readable hostname (e.g.\ \texttt{itcorp.com}) to an IPv4 address
before a socket can be opened.

\subsubsection*{Netualizer Stack Architecture}

The \texttt{TcpClient} project uses the following layer chain:

\begin{center}
  \texttt{phy \textrightarrow\ ethernet \textrightarrow\ ip \textrightarrow\ tcp \textrightarrow\ raw}
\end{center}

\begin{figure}[H]
  \centering
  \includegraphics[width=0.82\textwidth]{img_tcp_arch.png}
  \caption{Netualizer stack diagram for the TCP Client project (CLIENT node with raw/tcp/ip/ethernet/phy layers)}
  \label{fig:tcp_arch}
\end{figure}

\subsubsection*{Lua Script and Execution}

The automation script \texttt{TcpClient.lua} performs the following sequence:
(1)~wait for all stack layers to initialise, (2)~resolve
\texttt{itcorp.com} via the integrated DNS resolver, (3)~open a TCP socket to
port~80, (4)~send an HTTP~GET request, (5)~listen for incoming events.

\begin{figure}[H]
  \centering
  \includegraphics[width=0.82\textwidth]{img_tcp_code.png}
  \caption{TcpClient.lua source code with console output showing DNS resolution to \texttt{134.173.42.59} and TCP connection from \texttt{192.168.21.10}}
  \label{fig:tcp_code}
\end{figure}

\subsubsection*{Results and Analysis}

\begin{resultblock}{Observed Output}
\begin{itemize}[nosep]
  \item DNS resolution: \texttt{itcorp.com} $\rightarrow$ \texttt{134.173.42.59} \textbf{[OK]}
  \item Source IP: \texttt{192.168.21.10}
  \item TCP socket opened on port~80; connection packets generated
  \item Layer-3 routing table and Layer-4 handshake confirmed intact
\end{itemize}
\end{resultblock}

\textbf{Engineering Observation:} DNS resolution adds non-trivial latency
before the TCP handshake. In production IoT deployments, IP addresses are
typically hardcoded or resolved once at boot and cached, since DNS lookup
overhead (typically 20--200\,ms) is unacceptable for time-critical sensing
loops. For local node communication, mDNS (Section~\ref{sec:mdns}) eliminates
the external DNS dependency entirely.

% ── 3.2 MQTT ─────────────────────────────────────────────────────────────────
\subsection{Example 2: MQTT Sensor (Publish/Subscribe Telemetry)}
\label{sec:mqtt}

\subsubsection*{Theoretical Background}

\textbf{MQTT} (Message Queuing Telemetry Transport)~\cite{MQTT5} is a
lightweight broker-based publish/subscribe protocol engineered for
resource-constrained IoT devices. Key design properties:

\begin{itemize}[nosep]
  \item Runs over TCP; provides session persistence across reconnects
  \item Fixed header: only \textbf{2 bytes} (compared to HTTP's $\geq$200 bytes)
  \item Three Quality-of-Service tiers: QoS~0 (at most once), QoS~1
    (at least once), QoS~2 (exactly once)
  \item Hierarchical topic model: e.g.\ \texttt{home/living\_room/temperature}
  \item Broker decouples publishers and subscribers --- neither needs to know
    the other exists
\end{itemize}

This architecture enables a single broker to fan out one sensor's telemetry
to dozens of subscribers (dashboards, alerting systems, storage pipelines)
with no per-subscriber overhead on the sensor device.

\subsubsection*{Netualizer Stack Architecture}

\begin{center}
  \texttt{phy \textrightarrow\ ethernet \textrightarrow\ ip \textrightarrow\ tcp \textrightarrow\ mqtt}
\end{center}

\begin{figure}[H]
  \centering
  \includegraphics[width=0.82\textwidth]{img_mqtt_arch.png}
  \caption{Netualizer stack diagram for the MQTT Sensor project (SENSOR node with mqtt/tcp/ip/ethernet/phy layers)}
  \label{fig:mqtt_arch}
\end{figure}

\subsubsection*{Lua Script and Execution}

\texttt{MQTTSensor.lua} configures a static IP (\texttt{192.168.21.20})
to prevent gateway-drop faults, connects to a public sandbox broker at
\texttt{91.121.93.94}, and launches an event-driven publish loop that
continuously pushes \texttt{dataup~42} payloads.

\begin{figure}[H]
  \centering
  \includegraphics[width=0.82\textwidth]{img_mqtt_code.png}
  \caption{MQTTSensor.lua source code and Events panel showing the continuous \texttt{dataup~42} publish loop with per-event timestamps}
  \label{fig:mqtt_code}
\end{figure}

\subsubsection*{Results and Analysis}

\begin{resultblock}{Observed Output}
\begin{itemize}[nosep]
  \item Exit code: \texttt{done (0) <= MQTTSensor.lua}
  \item Event loop active; periodic \texttt{dataup 42} interrupts in Events panel
  \item Timestamps recorded: \texttt{Fri 06/05/2026 10:1x:xx} per publish cycle
  \item IP configuration: \texttt{Fup 192.168.21.20} confirmed in Events
\end{itemize}
\end{resultblock}

\textbf{Engineering Observation:} Static IP assignment was required because
the simulated gateway would drop packets from an unconfigured interface.
In production environments, a DHCP reservation (binding MAC to a fixed lease)
provides the same stability without manual address management. The continuous
publish loop with no connection renegotiation demonstrates MQTT's efficiency
for high-frequency telemetry --- a single persistent TCP connection serves
hundreds of publish operations.

% ── 3.3 CoAP ─────────────────────────────────────────────────────────────────
\subsection{Example 3: CoAP Client (Constrained Application Protocol)}
\label{sec:coap}

\subsubsection*{Theoretical Background}

\textbf{CoAP}~\cite{RFC7252} is a specialised RESTful protocol for
constrained nodes, standardised in RFC~7252. Its key differentiators from
HTTP:

\begin{center}
\renewcommand{\arraystretch}{1.2}
\begin{tabular}{lll}
\toprule
\textbf{Property} & \textbf{HTTP} & \textbf{CoAP} \\
\midrule
Transport     & TCP           & UDP (port 5683) \\
Header size   & $\geq$200 B   & 4~bytes base \\
Reliability   & TCP stream    & CON/NON messages \\
Model         & Client-Server & Client-Server (REST) \\
Multicast     & No            & Yes (UDP multicast) \\
\bottomrule
\end{tabular}
\end{center}

CoAP introduces \emph{Confirmable} (CON) messages for reliability: the
receiver must send an ACK, otherwise the sender retransmits with exponential
back-off. \emph{Non-confirmable} (NON) messages skip the ACK and are preferred
for periodic sensor readings where occasional loss is acceptable.

\subsubsection*{Netualizer Stack Architecture}

\begin{center}
  \texttt{phy \textrightarrow\ ethernet \textrightarrow\ ip \textrightarrow\ udp \textrightarrow\ coap}
\end{center}

\begin{figure}[H]
  \centering
  \includegraphics[width=0.82\textwidth]{img_coap_arch.png}
  \caption{Netualizer stack diagram for the CoAP Client project (IOT CLIENT node with coap/udp/ip/ethernet/phy layers)}
  \label{fig:coap_arch}
\end{figure}

\subsubsection*{Lua Script and Execution}

\texttt{CoAPClient.lua} prompts for the target sensor IP, forms a
\texttt{coap://} URI, and dispatches a binary CoAP GET packet over UDP to
the simulated temperature transducer at \texttt{192.168.21.30:5683}.

\begin{figure}[H]
  \centering
  \includegraphics[width=0.82\textwidth]{img_coap_code.png}
  \caption{CoAPClient.lua source code and output showing the GET request to \texttt{coap://192.168.21.30:5683/temperature}}
  \label{fig:coap_code}
\end{figure}

\subsubsection*{Results and Analysis}

\begin{resultblock}{Observed Output}
\begin{itemize}[nosep]
  \item URI formed: \texttt{coap://192.168.21.30:5683/temperature}
  \item Binary CoAP GET dispatched over UDP successfully
  \item No TCP handshake delay; connectionless dispatch confirmed
  \item No packet fragmentation observed across the virtual link
\end{itemize}
\end{resultblock}

\textbf{Engineering Observation:} The elimination of the TCP handshake
reduces connection setup time from $\sim$3 round-trips to 0. For
battery-powered sensors that wake, transmit, and sleep, this difference
directly translates to energy savings. The trade-off is that application-level
reliability (CON messages) must be explicitly managed, adding code complexity.

% ── 3.4 UDP ──────────────────────────────────────────────────────────────────
\subsection{Example 4: UDP Sender (Connectionless Datagram)}
\label{sec:udp}

\subsubsection*{Theoretical Background}

\textbf{UDP} provides a minimal stateless transport layer with an 8-byte
fixed header. Unlike TCP, UDP offers:
\begin{itemize}[nosep]
  \item No connection setup or teardown
  \item No delivery guarantees or ordering
  \item No retransmission or flow control
  \item Sub-millisecond dispatch latency
\end{itemize}

IoT use cases that benefit from UDP's low overhead include real-time GPS
tracking, industrial sensor streaming ($>$10\,Hz), video/audio over constrained
links, and multicast status beacons. Occasional packet loss is acceptable in
these scenarios, whereas the CPU and battery cost of TCP state machine
maintenance is not.

\subsubsection*{Netualizer Stack Architecture}

\begin{center}
  \texttt{phy \textrightarrow\ ethernet \textrightarrow\ ip \textrightarrow\ udp}
\end{center}

\begin{figure}[H]
  \centering
  \includegraphics[width=0.82\textwidth]{img_udp_arch.png}
  \caption{Netualizer stack diagram for the UDP Sender project (UDP node with udp/ip/ethernet/phy layers)}
  \label{fig:udp_arch}
\end{figure}

\subsubsection*{Lua Script and Execution}

\texttt{UdpSender.lua} binds to \texttt{192.168.21.30:5000}, sets the
destination to \texttt{192.168.21.5:6000}, and immediately dispatches the
raw payload \texttt{"raw data being sent!\textbackslash n"} without any
handshake.

\begin{figure}[H]
  \centering
  \includegraphics[width=0.82\textwidth]{img_udp_code.png}
  \caption{UdpSender.lua source code and output showing instant datagram dispatch and \texttt{done (0)} exit}
  \label{fig:udp_code}
\end{figure}

\subsubsection*{Results and Analysis}

\begin{resultblock}{Observed Output}
\begin{itemize}[nosep]
  \item Source socket bound: \texttt{192.168.21.30:5000}
  \item Payload sent immediately: \texttt{"raw data being sent!"}
  \item Exit code: \texttt{done (0) <= UdpSender.lua}
  \item No waiting state; sub-millisecond execution confirmed
\end{itemize}
\end{resultblock}

\textbf{Engineering Observation:} Comparing UDP's instant exit with TCP's
blocking handshake confirms the fundamental latency trade-off. In an
industrial telemetry application sampling at 50\,Hz, TCP's handshake overhead
alone would consume $\sim$150\,ms per connection (assuming 50\,ms RTT) ---
equivalent to 7--8 lost samples. UDP eliminates this entirely.

% ── 3.5 mDNS ─────────────────────────────────────────────────────────────────
\subsection{Example 5: mDNS Responder (Zero-Configuration Discovery)}
\label{sec:mdns}

\subsubsection*{Theoretical Background}

\textbf{Multicast DNS (mDNS)}~\cite{RFC6762} (RFC~6762) enables hostname
resolution within local networks without a centralised DNS server. It operates
as follows:
\begin{enumerate}[nosep]
  \item A device wanting to resolve \texttt{sensor01.local} sends a UDP
    multicast query to \texttt{224.0.0.251:5353}
  \item The device owning that name responds with its IP address
  \item No infrastructure (DHCP server, DNS server) is required
\end{enumerate}

mDNS is the foundation of Apple Bonjour and Linux Avahi. In IoT contexts, it
allows sensors, gateways, and edge servers to discover each other
automatically when joining a network --- critical for plug-and-play deployment
of large sensor fleets.

DNS-SD (DNS Service Discovery) extends mDNS with PTR/SRV/TXT records that
advertise service types (e.g.\ \texttt{\_sip.\_udp.local},
\texttt{\_mqtt.\_tcp.local}), allowing subscribers to find brokers or APIs
without hardcoded addresses.

\subsubsection*{Netualizer Stack Architecture}

\begin{center}
  \texttt{phy \textrightarrow\ ethernet \textrightarrow\ ip \textrightarrow\ udp \textrightarrow\ rtp \textrightarrow\ player}
\end{center}

\begin{figure}[H]
  \centering
  \includegraphics[width=0.82\textwidth]{img_mdns_arch.png}
  \caption{Netualizer stack diagram for the mDNS project (mDNS node with player/rtp/udp/ip/ethernet/phy layers)}
  \label{fig:mdns_arch}
\end{figure}

\subsubsection*{Lua Script and Execution}

\texttt{mDNS.lua} starts an mDNS responder daemon that binds to the
multicast port and broadcasts a PTR record for \texttt{\_sip.\_udp.local}.

\begin{figure}[H]
  \centering
  \includegraphics[width=0.82\textwidth]{img_mdns_code.png}
  \caption{mDNS.lua source code and output showing responder bound to \texttt{192.168.21.18:5353} and PTR query broadcast}
  \label{fig:mdns_code}
\end{figure}

\subsubsection*{Results and Analysis}

\begin{resultblock}{Observed Output}
\begin{itemize}[nosep]
  \item mDNS responder address: \texttt{192.168.21.18:5353}
  \item PTR query issued for \texttt{\_sip.\_udp.local}
  \item Exit code: \texttt{done (0) <= mDNS.lua}
  \item Local service record compiled and advertised successfully
\end{itemize}
\end{resultblock}

\textbf{Engineering Observation:} mDNS removes the need for manual IP
provisioning at device bring-up, which is operationally critical when deploying
hundreds of sensors. Without mDNS, each sensor requires either a DHCP
reservation or a hardcoded IP --- both create operational overhead. The PTR
record observed (\texttt{\_sip.\_udp.local}) also demonstrates DNS-SD
compatibility, allowing SIP-capable IoT devices to advertise themselves on
the local segment.

% ── 3.6 Data Logging ─────────────────────────────────────────────────────────
\subsection{Data Logging and Timestamp Verification}
\label{sec:logging}

\subsubsection*{Concept}

Temporal accuracy is a foundational requirement for IoT systems. Every
telemetry event must carry an immutable, high-resolution timestamp to enable:
\begin{itemize}[nosep]
  \item Latency measurement between publish and receive
  \item Ordered event replay for debugging
  \item Audit trails for regulatory compliance
  \item Correlation of readings across distributed nodes
\end{itemize}

\subsubsection*{Implementation and Results}

Netualizer's \emph{Events} panel automatically appends a monotonic timestamp
to every network state transition. No additional scripting is required. The
logging was observed during the MQTT session (Section~\ref{sec:mqtt}).

\begin{figure}[H]
  \centering
  \includegraphics[width=0.90\textwidth]{img_logging.png}
  \caption{Netualizer Events panel during MQTT session: each \texttt{dataup~42} event carries a unique timestamp (format: \texttt{Fri 06/05/2026 10:1x:xx}), confirming ordered, non-repeating audit trail}
  \label{fig:logging}
\end{figure}

\begin{resultblock}{Observed Logging Behaviour}
\begin{itemize}[nosep]
  \item Timestamp format: \texttt{Fri 06/05/2026 10:1x:xx} (second-level resolution)
  \item All \texttt{dataup~42} MQTT events recorded with unique, strictly
    increasing timestamps
  \item Event sequence preserved in insertion order; no out-of-order entries
  \item IP configuration events (\texttt{Fup}) also logged with timestamps
\end{itemize}
\end{resultblock}

\textbf{Engineering Observation:} For production pipelines these logs should
be forwarded to a persistent indexed store (e.g.\ Quickwit) with structured
fields: \texttt{ts} (timestamp), \texttt{device\_id}, \texttt{sensor\_type},
\texttt{value}. This enables SQL-style queries (via Trino or DataFusion) and
real-time Grafana dashboards over the telemetry stream.

% ── 3.7 Failure Injection ────────────────────────────────────────────────────
\subsection{Network Failure Injection and Exception Analysis}
\label{sec:failure}

Robust IoT network design requires systematic fault injection to validate
that the protocol stack detects, reports, and fails gracefully under
misconfiguration, missing hardware, and environmental anomalies.
Three distinct failure scenarios were deliberately induced.

\subsubsection*{Scenario 1 --- Network-Layer Gateway Omission}

\textbf{Setup:} The IP layer was initialised without assigning a default
gateway, then a transmission was attempted.

\begin{figure}[H]
  \centering
  \includegraphics[width=0.82\textwidth]{img_fail1.png}
  \caption{Failure Scenario 1: Netualizer dialog reporting \textit{``Interface `ip' failed to start [error: gateway not set!]''} when the IP layer has no configured default gateway}
  \label{fig:fail1}
\end{figure}

\begin{alertblock}{Error Produced}
\texttt{Interface 'ip' failed to start [error: gateway not set!]}
\end{alertblock}

\textbf{Analysis:} The network layer correctly blocks all upward packet flow
when it cannot route. This mirrors real-world behaviour: a sensor that boots
without a valid gateway will discard all outbound packets silently at the IP
layer. Netualizer surfaces this as an explicit startup exception, allowing
engineers to catch misconfiguration at integration time rather than during
deployment.

\subsubsection*{Scenarios 2 and 3 --- Agent Absence and LoRa Hardware Fault}

\textbf{Setup:} A LoRa-based project was launched (a)~without any agent
connected to the controller, and (b)~without attaching the virtual RYLR896
transceiver module.

\begin{figure}[H]
  \centering
  \includegraphics[width=0.82\textwidth]{img_fail2_3.png}
  \caption{Failure Scenarios 2 \& 3: (top) \textit{``No agents were found''} when the controller has no connected agent; (bottom) \textit{``Interface RYLR896+LoRa+1 is not physically connected to this agent!''} when the LoRa transceiver is absent}
  \label{fig:fail23}
\end{figure}

\begin{alertblock}{Errors Produced}
\begin{enumerate}[nosep]
  \item \texttt{No agents were found.}
  \item \texttt{Interface RYLR896+LoRa+1 is not physically connected to this agent!}
\end{enumerate}
\end{alertblock}

\textbf{Analysis:}
\begin{itemize}
  \item \textbf{No agents found:} The controller-agent architecture of
    Netualizer requires at least one agent process to be running. This fault
    confirms that the simulation correctly enforces the presence of a virtual
    execution host before allowing any device configuration.
  \item \textbf{LoRa transceiver absent:} Physical-layer hardware is a
    mandatory prerequisite. Without a radio module attached, the LoRa MAC
    and upper layers cannot initialise. This directly models real IoT
    deployments where a missing antenna or disconnected RF module causes
    total communication failure --- a common field fault that is difficult
    to diagnose without structured error reporting.
\end{itemize}

% ══════════════════════════════════════════════════════════════════════════════
\section{Comparative Protocol Analysis}
\label{sec:comparison}

Table~\ref{tab:proto} synthesises the engineering trade-offs observed across
all five simulation runs. No single protocol dominates; the optimal choice is
always context-dependent.

\begin{table}[H]
\centering
\caption{IoT protocol comparison matrix derived from simulation results}
\label{tab:proto}
\renewcommand{\arraystretch}{1.35}
\begin{tabular}{>{\bfseries}p{1.6cm} p{1.4cm} p{2.1cm} p{1.8cm} p{4.4cm}}
\toprule
Protocol & Transport & Connection model & Header overhead & Primary IoT use case \\
\midrule
TCP   & TCP (L4) & Connection-oriented (3-way HS) & High ($\geq$20\,B)    & FOTA, REST APIs \\
UDP   & UDP (L4) & Connectionless (datagram)       & Low (fixed 8\,B)      & Real-time telemetry, GPS, streaming \\
MQTT  & TCP (L4) & Broker pub/sub (async)          & Very low (2\,B fixed) & Cloud telemetry, smart grids \\
CoAP  & UDP (L4) & Client-server REST (req/resp)   & Very low (4\,B base)  & Constrained mesh, battery sensors \\
mDNS  & UDP (L4) & Multicast broadcast (zero-cfg)  & Low--moderate         & Smart-home discovery, ad-hoc LANs \\
\bottomrule
\end{tabular}
\end{table}

\textbf{Key design insight:} UDP-based protocols (CoAP, mDNS, raw UDP)
minimise overhead at the cost of application-level reliability management.
TCP-based protocols (direct TCP, MQTT) provide guaranteed delivery but consume
more energy and introduce higher latency. The choice must be driven by the
specific application's requirements for reliability, energy budget, and
data rate.

% ══════════════════════════════════════════════════════════════════════════════
\section{Security Considerations}
\label{sec:security}

None of the five implemented examples includes transport-layer or
application-layer security. This is a deliberate scope decision for a
foundational Track~B submission. However, deploying these protocols in
production without security hardening exposes the system to serious threats:

\begin{table}[H]
\centering
\caption{Threat model and mitigations per protocol}
\label{tab:security}
\renewcommand{\arraystretch}{1.3}
\begin{tabular}{p{1.4cm} p{3.5cm} p{5.5cm}}
\toprule
\textbf{Protocol} & \textbf{Primary threat} & \textbf{Recommended mitigation} \\
\midrule
MQTT  & Unauthenticated broker access; topic eavesdropping & TLS~1.2+ channel; X.509 client certificates; topic-level ACLs \\
CoAP  & Replay attacks; UDP amplification DDoS & DTLS with PSK or certificates; frame-counter validation (RFC~8613) \\
mDNS  & Spoofed PTR/A records; local traffic interception & Network segmentation (VLAN); DNSSEC where supported \\
TCP   & Plaintext payload interception; man-in-the-middle & TLS~1.3 end-to-end \\
UDP   & Payload spoofing; no integrity check & DTLS or application-layer HMAC \\
\bottomrule
\end{tabular}
\end{table}

A natural extension of this project would implement \textbf{HMAC-SHA256} data
integrity verification on the MQTT payload and demonstrate a replay-attack
simulation (reinjecting a captured \texttt{dataup} packet) to qualify for the
Security Bonus criterion.

% ══════════════════════════════════════════════════════════════════════════════
\section{Conclusion}
\label{sec:conclusion}

Seven simulation scenarios were successfully executed within Netualizer,
covering five IoT protocol stacks and two supplementary engineering scenarios.
The following conclusions are drawn:

\begin{enumerate}
  \item \textbf{TCP and MQTT} deliver reliable, ordered communication but
    impose connection overhead and higher energy consumption --- suitable for
    non-time-critical, high-reliability applications such as FOTA and
    centralised telemetry.

  \item \textbf{UDP and CoAP} minimise handshake latency and header overhead,
    making them preferable for battery-limited, high-frequency edge devices.
    CoAP adds optional application-level reliability (CON messages) without
    reverting to TCP.

  \item \textbf{mDNS} enables automatic service discovery with zero
    infrastructure dependency, which is operationally essential for deploying
    and scaling IoT fleets without per-device manual configuration.

  \item \textbf{Netualizer} correctly enforces multi-layer dependencies,
    produces actionable fault diagnostics, and provides high-resolution
    timestamps for all state transitions --- making it a valuable tool for
    protocol validation before physical deployment.

  \item \textbf{Security} is absent from all base implementations and must
    be addressed before any production deployment, with DTLS and HMAC as the
    highest-priority additions for UDP-based protocols.
\end{enumerate}

The comparative protocol matrix (Table~\ref{tab:proto}) and the threat model
(Table~\ref{tab:security}) provide practical references for IoT protocol
selection and security planning in constrained deployment scenarios.

% ══════════════════════════════════════════════════════════════════════════════
\section*{Reproducibility}
\label{sec:repro}

All scenarios run in Netualizer with the virtual network \texttt{192.168.21.0/24}.
To reproduce:
\begin{enumerate}[nosep]
  \item Install Netualizer (Controller + Agent) from \url{https://www.l7tr.com/site/netualizer}
  \item Open the project directory for each example (e.g.\ \texttt{TcpClient/})
  \item Apply static IP settings where noted (MQTT: \texttt{192.168.21.20})
  \item Run the corresponding \texttt{*.lua} script
  \item Observe the Output and Events panels
\end{enumerate}
No external hardware or additional dependencies are required.

% ══════════════════════════════════════════════════════════════════════════════
\begin{thebibliography}{9}

\bibitem{RFC7228}
C.\ Bormann, M.\ Ersue, A.\ Keranen,
\textit{Terminology for Constrained-Node Networks},
RFC~7228, IETF, May 2014.
\url{https://www.rfc-editor.org/rfc/rfc7228}

\bibitem{MQTT5}
OASIS Standard,
\textit{MQTT Version 5.0},
OASIS, March 2019.
\url{https://docs.oasis-open.org/mqtt/mqtt/v5.0/mqtt-v5.0.html}

\bibitem{RFC7252}
Z.\ Shelby, K.\ Hartke, C.\ Bormann,
\textit{The Constrained Application Protocol (CoAP)},
RFC~7252, IETF, June 2014.
\url{https://www.rfc-editor.org/rfc/rfc7252}

\bibitem{RFC6762}
S.\ Cheshire, M.\ Krochmal,
\textit{Multicast DNS},
RFC~6762, IETF, February 2013.
\url{https://www.rfc-editor.org/rfc/rfc6762}

\bibitem{RFC8613}
O.\ Bergmann et al.,
\textit{Object Security for Constrained RESTful Environments (OSCORE)},
RFC~8613, IETF, July 2019.
\url{https://www.rfc-editor.org/rfc/rfc8613}

\bibitem{Herrero2023}
Herrero,
\textit{IoT Architecture and Networking: Architecture and Protocols}, 2023.

\bibitem{Edwards2024}
Edwards,
\textit{IoT Security: Threat Models and Controls}, 2024.

\bibitem{NISTIR8228}
NIST,
\textit{NISTIR 8228: Considerations for Managing IoT Cybersecurity and Privacy Risks}, 2019.
\url{https://doi.org/10.6028/NIST.IR.8228}

\end{thebibliography}

\end{document}
